% Generated draft: requires mathematical and source review.
\documentclass[11pt]{article}
\usepackage[margin=25mm]{geometry}
\usepackage{fontspec}
\setmainfont{Latin Modern Roman}
\usepackage{amsmath,amssymb,mathtools,unicode-math}
\setmathfont{Latin Modern Math}
\usepackage{enumitem,needspace,xurl,hyperref}
\setlength{\parindent}{0pt}
\setlength{\parskip}{5pt}
\setlength{\emergencystretch}{3em}
\newcommand{\N}{\mathbb{N}}
\newcommand{\Z}{\mathbb{Z}}
\newcommand{\Q}{\mathbb{Q}}
\newcommand{\R}{\mathbb{R}}
\newcommand{\C}{\mathbb{C}}
\newcommand{\F}{\mathbb{F}}
\newcommand{\A}{\mathbb{A}}
\newcommand{\PP}{\mathbb{P}}
\newcommand{\E}{\mathbb{E}}
\newcommand{\eps}{\varepsilon}
\newcommand{\dd}{\,\mathrm{d}}
\begin{document}
\section*{UnsolvedMath ID: 701549}
\textbf{Source problem code:} catalogue entry
\subsection{Definitions and mathematical statement}
Let \(G\) be a group with identity element \(1_G\). For \(S\subseteq G\), let \(\langle S\rangle\) be the subgroup generated by \(S\), and write \(\langle g,S\rangle=\langle\{g\}\cup S\rangle\). An element \(g\in G\) is a \emph{non-near generator} if
\[
\forall S\subseteq G,\qquad [G:\langle g,S\rangle]<\infty\ \Longrightarrow\ [G:\langle S\rangle]<\infty.
\]
Define the \emph{lower near Frattini subgroup} by
\[
\lambda(G)=\{g\in G:g\text{ is a non-near generator}\}.
\]
A subgroup \(M\leq G\) is \emph{nearly maximal} if
\[
[G:M]=\infty\quad\text{and}\quad
\forall H\leq G,\qquad M\subsetneq H\ \Longrightarrow\ [G:H]<\infty.
\]
The \emph{upper near Frattini subgroup} is
\[
\mu(G)=\bigcap_{\substack{M\leq G\\M\text{ nearly maximal}}}M.
\]
Here an intersection over an empty family of subgroups is taken to be \(G\); the supplied record does not explicitly state this convention. The source states that \(\lambda(G)\) and \(\mu(G)\) are characteristic subgroups and that \(\lambda(G)\leq\mu(G)\). When they coincide, their common value is called the \emph{near Frattini subgroup} and is denoted by \(\psi(G)\).

The two questions proposed by M. K. Azarian are:
\begin{enumerate}
\item[(a)] For every group \(G\) that is the knot group of a product of knots, is
\[
\lambda(G)=\mu(G)=\{1_G\}?
\]
\item[(b)] For every cable knot group \(G\), is
\[
\lambda(G)=\mu(G)=\{1_G\}?
\]
\end{enumerate}
These equalities express the supplied questions \(\psi(G)=1\) while explicitly including the equality needed for \(\psi(G)\) to be defined.

The knot terminology requires qualification. Under the customary interpretation of knots in \(S^3\), the knot group of \(K\) is \(\pi_1(S^3\setminus K)\), considered up to isomorphism. The supplied record does not specify the ambient manifold or knot category, define a product of knots, or state the permitted cabling parameters. Interpreting a product as a finite connected sum is a possible reading, but is not imposed here. Thus the exact knot classes quantified over in (a) and (b) remain subject to these unresolved source qualifications.
\subsection{Short English statement}
Azarian asks whether the lower and upper near Frattini subgroups are both trivial for every knot group arising from a product of knots, and for every cable knot group. This formulation includes the equality required to define \(\psi(G)\). The precise knot classes, product operation, and cabling conventions require clarification from the source.
\subsection{Sources}
\begin{itemize}
\item Supplied local UnsolvedMath registry record: permanent ID \(701549\), exact source problem code \texttt{catalogue entry}, title \emph{Kourovka Notebook Problem 19.1}; upstream revision \texttt{d371ee917ac9647fbb06cd6fc529deec6f3dda82}. The record was created and updated on 2026-10-16.
\item The supplied attribution names M. K. Azarian as proposer and gives 2018 as the proposed year.
\item The supplied background cites \emph{The Kourovka Notebook (Unsolved Problems in Group Theory)}, 19th issue (2018), original problem 19.1, and the 21st issue edition (2026), identified as arXiv:1401.0300v48. The supplied collection description names E. I. Khukhro and V. D. Mazurov as editors.
\item The local record labels the problem \emph{open}. This is supplied source metadata, not an independently verified mathematical or literature assessment. The cited notebook editions were not independently consulted for this draft.
\end{itemize}
\end{document}
